\chapter{知识的共享}
\label{chap:ket-sharing}

一张花卉照片可以用于识别物种，也可以用于判断病害和确定花朵位置。三个模型分别训练时，都可能需要从照片中学习边缘、纹理和形状；让它们共用这些计算，便有机会减少重复估计。不过，物种分类不必保留的细小斑点，可能正是病害判断的依据；分类时可以忽略的绝对位置，定位时却必须输出。我们要解决的不是怎样让三个模型处处相同，而是怎样利用共同规律，同时给每项预测保留所需的差异。

让不同预测要求共同参与训练，使一项任务的经验能够改变另一项任务所用的模型，这称为\term{多任务学习}（multi-task learning, MTL）。本章以常驻的花卉类别、病害判断和对象定位任务$\mathcal T_1,\mathcal T_2,\mathcal T_3$为例；推广时以$K$表示参与任务数，这个例子中$K=3$。要判断共享是否有益，先为每个任务约定怎样计量错误。任务$k$的预测器为$f_k$，损失$\ell_k$对一次预测计分，联合分布$\mathcal P_k$描述该任务未来面对的输入与标签。对这个分布取平均，得到需要尽量降低的任务风险：
\begin{equation}
  R_k(f_k)=
  \mathbb E_{(\bm{x},y_k)\sim\mathcal P_k}
  [\ell_k(f_k(\bm{x}),y_k)].
  \label{eq:ket-sharing-task-risk}
\end{equation}
标签$y_k$分别表示类别、病害状态或位置，输出空间和损失单位都可以随任务改变。对训练集$S_{\mathrm{train},k}$取样本平均所得的经验风险记为$\hat R_k$，对当前训练批次计算的损失记为$L_k$。降低$L_k$说明模型更符合本批监督，还不能证明共享降低了未来风险；后者需要用未参与拟合和选择的评价数据估计。

共享方案需要同时说明谁与谁学习、共同使用什么、怎样取舍各项错误，以及为此投入多少标签和计算。本章的四个主体节分别处理这四项设计，比较一项设计时尽量固定其余条件。它们可以组合，也会相互制约，并不是四种互斥算法。第~\ref{chap:ket-transfer-adaptation}章让已有经验服务于新目标，本章则让多个当前任务一起形成经验；如果模型已经分别学好，只需再组合它们的计算或融合参数，应回到第~\ref{chap:ket-integration}章。

区分共享方式，还要看信息在哪个阶段发生联系。共用图像并不等于共用表示，共用表示也不等于输出必须满足同一关系：信息可以在数据组织、前向计算或训练约束中交换。训练时又须分清目标、估计和更新这三个环节。任务权重规定要降低哪些错误，采样与缺失标签处理决定实际取得什么监督，梯度协调再把这些监督变成参数变化。即使最后一步安排得当，前面漏掉的任务或偏向某类样本的标签也不会自动得到补偿。

\section{共享任务的组织}
\label{sec:ket-sharing-organization}

类别、病害和定位都读取照片，但部署责任可能不同。若三项输出都要交付，就必须分别检查；若最终只交付类别，其他预测可以仅在训练中提供帮助。先明确这种收益要求，再测量任务之间的影响，才能决定是否全部共同训练或只让部分任务成组。主辅地位回答保护谁的收益，分组与层级回答谁和谁共享，两者可以同时描述一个方案。试训若表明病害持续受损，分组也应随证据修订。

\subsection{候选任务与收益要求的确定}
\label{sec:ket-sharing-candidates}

假设终端必须同时给出类别、病害和位置，就需要为三项输出各设最低标准。病害漏检的成本与定位偏差的成本不相同，不能先设三个相等权重，再把所得结果当作部署要求。除绝对标准外，还可以限制共享相对分别训练的损失：在相同评价条件和约定预算下，以$R_k^{\mathrm{single}}$为单任务基线风险，以$R_k^{\mathrm{joint}}$为共同训练风险。沿用式~\eqref{eq:ket-reuse-risk-gain}中正收益表示改善的约定，对必须交付的任务集合$C$施加
\begin{equation}
  \Delta_k=R_k^{\mathrm{single}}-R_k^{\mathrm{joint}}
  \geq-\epsilon_k,
  \qquad k\in C,
  \label{eq:ket-sharing-benefit-constraint}
\end{equation}
其中$\epsilon_k\geq0$是事先允许的风险增加量。共享降低风险时$\Delta_k>0$，增加风险时$\Delta_k<0$；取$\epsilon_k=0$便要求不退化。这项相对约束不能代替绝对最低标准：两个方案都不合格时，相对收益仍可能为正。不同损失单位下的$\Delta_k$也不能直接相加，跨指标汇总见第~\ref{sec:ket-evaluation-sharing}节。由于风险由有限数据估计，验收还须处理误差；式中写的是要求，不是训练算法已经提供的保证。

若只需部署类别识别，花瓣颜色或叶缘形状预测可以作为\term{辅助任务}（auxiliary task）。这类任务给训练增加监督，但价值由主任务是否受益决定，不要求它自己成为独立交付能力。属性监督可能迫使共同表示保留颜色或叶缘信息，使类别头不必仅凭有限物种标签学习所有区分线索。Caruana用任务共用隐藏表示、分别使用输出单元的方式说明了这条影响路径\scite{caruana1997}。

辅助监督保留的信息也可能与主任务所需信息不同。例如，跨季节识别物种需要容忍某些色差，而一个细分光照色差的辅助目标可能占用共同表示的有限容量。判断是否保留该目标，应看重复训练中的类别收益及其波动，再核算标签费用和部署费用，不能以属性损失下降替代检查。辅助头若只在训练时提供损失，部署时可删除；若类别预测实际读取属性输出，删除它便会改变主任务计算，必须连同这条依赖重新检验。

辅助预测问题还可以从已有数据构造。Ando和Zhang让多个预测任务共同学习结构，同时允许任务特有函数分别拟合目标；共同结构限制的是各任务可选择的表示或假设空间\scite{ket-ando2005-structure}。构造目标是否有效、标签是否可靠及是否泄漏待预测答案，由第~\ref{part:efficient-knowledge-use}部分讨论。这里先把候选目标与收益标准对应起来，下一步再测量它在共同训练中实际帮助谁。

\subsection{任务间共享关系的估计}
\label{sec:ket-sharing-affinity}

物种与病害标签相关，可能因为物种确实影响病害，也可能因为采样时某些物种只拍到了患病植株。仅凭这种相关性，无法知道类别监督会不会把病害模型引向不可靠的线索。我们需要在已经选定的结构中测量训练影响：从同一个参数状态出发，单独采用某项任务的更新，其他任务的损失怎样变化？

把阶段$t$的共享参数记为$\bm{\theta}_{\mathrm{s}}^{(t)}$，任务私有参数记为$\bm{\theta}_k^{(t)}$。为隔离共享块的作用，固定当前批次和所有私有参数，先计算任务$k$的共享梯度$\bm g_k^{(t)}$，再用共同步长$\eta>0$试做一步：
\begin{align}
  \bm g_k^{(t)}
    &=\nabla_{\bm{\theta}_{\mathrm{s}}}L_k^{(t)},
  \\
  \widetilde{\bm{\theta}}_{\mathrm{s},k}^{(t)}
    &=\bm{\theta}_{\mathrm{s}}^{(t)}-\eta\bm g_k^{(t)}.
  \label{eq:ket-sharing-virtual-step}
\end{align}
这一步只用于测量，完成后恢复原参数，因此称为虚拟更新；测量下一个任务时也从原状态开始。记$L_{j\mid k}^{(t)}$为试用任务$k$的共享块更新后，任务$j$在固定批次上的损失。为将变化与任务$j$自身的当前损失相比，在$L_j^{(t)}>0$时定义有向的\term{任务亲和性}（task affinity）：
\begin{equation}
  A_{k\to j}^{(t)}
  =1-\frac{L_{j\mid k}^{(t)}}{L_j^{(t)}}.
  \label{eq:ket-sharing-affinity}
\end{equation}
这个数表示任务$k$的一步更新使任务$j$的损失相对原值降低多少：正值是降低，负值是升高。Fifty等把训练过程中测得的这类影响汇总起来，作为搜索任务组的依据\scite{ket-fifty2021-affinity}。

一次测量仍受批次与训练阶段影响。比较不同方向时应固定步长、共享层范围和对应的评价批次，并在多个批次、阶段记录平均值及波动。损失接近零时，相对变化可能被很小的分母放大，应补报绝对变化，或明确另用什么归一化尺度。训练批次给出的只是当前训练影响；若用验证集选择分组，该验证集已经参与模型选择，最终测试须另外保留。

为什么梯度内积可以提示相互影响？固定私有参数后，对任务$j$的损失沿任务$k$的共享更新作泰勒展开；共享梯度局部Lipschitz连续使余项为二阶量，从而有
\begin{equation}
  L_{j\mid k}^{(t)}-L_j^{(t)}
  =-\eta(\bm g_j^{(t)})^{\top}\bm g_k^{(t)}
  +O(\eta^2).
  \label{eq:ket-sharing-local-effect}
\end{equation}
在非零内积且步长足够小时，一阶项决定变化方向：正内积降低$j$的损失，负内积提高它。原始内积交换$j,k$后不变，但有限步长的结果还受两项损失各自的曲率影响，亲和性又以不同原损失作分母。因此，内积的对称性不能推出测得的亲和性也对称。

用一个可手算的教学构造检查这一区别。设共享标量$\theta$表示一种特征强度，类别损失为$L_{\text{类}}(\theta)=\tfrac12(\theta-1)^2$，属性损失为$L_{\text{属}}(\theta)=5(\theta-0.1)^2$。在零点两梯度都为$-1$，取$\eta=0.3$时任一任务提出的参数都是$0.3$。此时类别损失由$0.5$降为$0.245$，属性损失由$0.05$升为$0.2$，故两个方向分别得到$A_{\text{属}\to\text{类}}=0.51$和$A_{\text{类}\to\text{属}}=-3$。属性损失曲率较大，这一步跨过其最优点并走得过远，才造成单向帮助。若把步长减为$0.01$，两项损失都会降低；不对称并非梯度内积本身造成。

这个标量例子还可以区分相互帮助与相互干扰。两项正曲率二次损失的最优点都位于当前位置右侧时，足够小的正向移动可同时降低损失。若把病害损失设为$\tfrac12(\theta+1)^2$，其零点梯度为$1$，便与类别梯度$-1$要求相反的移动。这些是假设损失的局部计算，不是对真实花卉任务的关系认定；更换表示、标签或共享层，或者进入另一个训练阶段，都要重新测量。

这里的测量与第~\ref{sec:ket-reuse-task-relations}节的顺序迁移不同。Taskonomy先固定来源表示，再训练目标读出以考察表示可复用性\scite{ket-zamir2018-taskonomy}。共同训练则允许目标损失反过来改变共享表示，所以能从类别表示读出属性，不代表加入属性训练一定帮助类别。只部署类别时应着重检查属性对类别的方向；三项共同交付时则要逐项记录帮助、受损和估计不确定性。

\subsection{任务分组与共享层级的确定}
\label{sec:ket-sharing-grouping}

假设当前骨干中类别与定位的测量结果较相容，病害却持续要求另一组纹理，可以把前两项成组、病害单独保留，再与全部共享和全部分别训练比较。亲和性用来提出候选，不直接确定最终分组：三个任务同训会改变梯度合成、容量竞争和各项采样机会，两两帮助并不足以保证整组帮助。

分组比较还要把模型大小算进去。将一个网络拆成三个同样大的网络会增加保存参数；将三个小网络替换成一个大骨干，也不只改变了共享关系。Standley等在不同网络大小、数据规模和计算预算下研究任务分组，说明关系证据须连同这些条件解释\scite{ket-standley2020-grouping}。可先用局部测量筛选少量分组，再在相同总预算下配置并完整训练各组，按逐任务验证指标选方案，最后在独立测试集评价。候选筛选省去的是一部分试训，并不能代替整组验证。

组织单位也可能不是不同预测要求，而是同一要求在不同地区的副本。为说明层级共享，暂时只保留类别识别，用$a$索引植物园，$c(a)$表示其地区组。全体副本可共同学习边缘，地区组与植物园各自保留拍摄条件或校准差异。只有对同维且坐标对应的参数块，才可用加法表示全体、组和个体的贡献：
\begin{equation}
  \bm{\theta}_a
  =\bm{\mu}+\bm b_{c(a)}+\bm\delta_a,
  \label{eq:ket-sharing-hierarchy}
\end{equation}
共同中心$\bm{\mu}$由全体使用，$\bm b_{c(a)}$由同组植物园使用，$\bm\delta_a$仅属于植物园$a$。这只是结构分解，尚未指定估计方法；若无进一步约束，各分量还可能有多种等价分解。地区结构提供共享候选而非正确性保证，相邻植物园采用不同相机时，地理分组未必反映观测差异。

植物园不能汇集数据时，可以各保留模型而协调训练。联邦多任务学习MOCHA以节点模型$\bm w_a$和联系节点的关系矩阵$\bm\Omega$联合组织经验损失与关系正则，并交替处理参数和关系结构\scite{ket-smith2017-mocha}。关系固定后，各节点利用本地数据近似求解子问题，再参与全局协调，不必被限制成同一个预测器。其优化分析对凸损失、正则项、局部求解误差和节点掉线都有条件，不能直接承诺任意非凸网络的行为。数据不离开节点也不保证更新没有信息泄漏；隐私条件见\BookChapterRef{chap:trust-privacy}{卷3《基础、理论与可信性》第16章“隐私保护”}，通信与执行方式由\BookPartRef{part:systems}{卷12《系统》第1部分}展开。

任务单位需要在异构或指令数据中继续明确。同一预测要求改写十个模板，不自动成为十项独立能力；同一模板换成另一种语言，也不抹去语言与数据差异。分别记录语义、来源、语言和模板，才能决定分组及评价的汇总单位。确定哪些单位可以交换经验后，还须决定这种交换发生在共同计算、参数关系还是输出关系中。

\section{可共享知识的设计}
\label{sec:ket-sharing-structure}

同组任务未必需要处处使用相同网络。类别与病害可以共同提取特征，也可以各自计算、只让参数估计受到共同先验约束；定位结果还可以作为病害预测的输入，或者只在训练中限制病害位置。下面按共同计算、参数先验和输出关系三个位置展开。表示若由共同参数计算，二者就是同一共享路径的不同描述，不应另当作互不相干的两类方法。

让信息共同使用以前，必须确认接口所指对象相同。图像裁剪、坐标原点和标签对应错位时，同形张量也可能无法交换。T5用任务前缀和文本答案把不同要求接入文本到文本接口；这统一了形式，却不统一任务难度、数据规模，也不意味着其全部实验均采用同时联合训练\scite{raffel2020t5}。MultiModel为图像、语音和文本保留各自输入输出模块，再让它们使用共同内部计算\scite{ket-kaiser2017-multimodel}。两种接口安排都可能支持共享，应检查的是它们是否保留当前任务需要的信息。

\subsection{共同计算结构的设计}
\label{sec:ket-sharing-computation}

若多个任务共用一段映射，各任务的损失就都可能改变它。此时要确定共享到哪一层、哪些分支交换特征，以及连接或混合比例是否随请求变化。固定连接可以学习其系数，条件连接则在使用时按任务或输入选择；只依赖任务身份的选择对同一任务仍然固定。下面以固定连接和逐输入混合为主要对照，两者都允许保留私有分支。

\subsubsection{固定共享范围的设计}
\label{sec:ket-sharing-fixed}

对同一张花卉照片，可以只计算一次共同特征，再由不同预测部分分别读出类别、病害和位置。将共同映射称为骨干，将各项输出映射称为任务头，其计算为
\begin{equation}
  \bm z=h_{\bm{\theta}_{\mathrm{s}}}(\bm x),
  \qquad
  \hat y_k=q_{\bm{\theta}_k}(\bm z).
  \label{eq:ket-sharing-backbone}
\end{equation}
这里$\bm z$是共同表示，$h_{\bm{\theta}_{\mathrm{s}}}$由共享参数控制，$q_{\bm{\theta}_k}$用任务私有参数把表示转成$\hat y_k$。各损失经任务头回传到同一个骨干，使其学习受多项要求限制。只共享浅层时，深层分支可以分别保留空间细节；共享更多层则减少独立计算，却要求同一高层表示保留所有必要信息。分类可使用位置不敏感的汇总，定位必须保留空间对应，这个差异解释了为何共享深度要靠具体任务检验。

共同表示不要求各任务具有同一预测函数。为看清这一点，暂用同一$d$维输入上的标量线性预测作简化：$\bm U\in\mathbb R^{d\times r}$给出$r$个共同方向，任务$k$用系数$\bm a_k\in\mathbb R^r$选择组合。将各任务参数并列成矩阵，便有
\begin{equation}
  \hat y_k=\bm a_k^{\top}\bm U^{\top}\bm x,
  \qquad
  \bm W=\bm U\bm A,
  \label{eq:ket-sharing-subspace}
\end{equation}
其中$\bm A\in\mathbb R^{r\times K}$的第$k$列为$\bm a_k$，$\bm W\in\mathbb R^{d\times K}$的第$k$列为该任务参数。列向量可以不同，但都只能落在$\bm U$张成的子空间内。多任务特征学习通过共同基和系数行稀疏约束，鼓励任务使用较少的一组特征；它与上述预定$r$的表示有关，却不是同一约束\scite{ket-argyriou2006-features}。训练联合估计共同结构与任务系数，部署按选定映射计算预测。限制太强会排除任务所需方向，限制太弱则可能不再有效减少估计自由度。

若完全共用骨干限制过强，可以保留分支，只在指定层交换激活。Cross-stitch网络用可训练系数线性混合形状已经对齐的分支特征。以同一位置、通道的两个标量激活$h_1,h_2$为例，混合后供两个分支使用的值为
\begin{equation}
  \begin{bmatrix}
    \widetilde h_1\\ \widetilde h_2
  \end{bmatrix}
  =
  \begin{bmatrix}
    a_{1,1}&a_{1,2}\\a_{2,1}&a_{2,2}
  \end{bmatrix}
  \begin{bmatrix}
    h_1\\h_2
  \end{bmatrix}.
  \label{eq:ket-sharing-cross-stitch}
\end{equation}
矩阵的对角项缩放本分支激活，非对角项把另一分支激活引入当前分支。系数与网络参数一起学习，但不是概率，行和也无需为1。上式展开的是两个分支的一个位置；更多任务可在事先选定的分支间设置对应混合\scite{ket-misra2016-cross-stitch}。

假设教学系数的两行为$(0.9,0.1)$和$(0.2,0.8)$，输入激活为$2,6$，则输出为$2.4,5.2$。反向传播同时更新系数与原分支参数。部署时这组系数保持训练后的值，不因照片改变而由另一个网络重算；照片改变的是被混合的激活。连接位置和参与分支也预先固定，这正是它与输入条件门控的区别。

上述三种固定结构都在共同估计与任务自由度之间取舍。共享少，可能重复学习纹理；共享多，可能没有足够方向保留病害和位置。比较共享层数或交换位置时要控制参数、特征维数和计算预算，否则更宽分支带来的收益会混入共享机制的效果。

\subsubsection{条件共享范围的设计}
\label{sec:ket-sharing-conditional}

固定混合比例还可能不足以处理同一任务内部的变化。清晰叶片上的病害判断可以利用斑点，遮挡输入却可能需要形状和上下文。我们可以保留多种映射，让当前输入决定它们各占多少，而不为每种情形另建完整模型。

\term{多门控混合专家模型}（Multi-gate Mixture-of-Experts, MMoE）把这些可复用映射称为专家，并为每项任务设置一个计算混合比例的门控网络。设$E$个专家$\bm e_m(\bm x)$输出同维向量，门控分数$a_{k,m}(\bm x)$由任务$k$的网络产生。用softmax保证权重非负且和为1，再对专家向量作加权和：
\begin{align}
  g_{k,m}(\bm x)
    &=\frac{\exp(a_{k,m}(\bm x))}
      {\sum_{r=1}^{E}\exp(a_{k,r}(\bm x))},
  \\
  \bm z_k(\bm x)
    &=\sum_{m=1}^{E}g_{k,m}(\bm x)\bm e_m(\bm x),
  \nonumber\\
  \hat y_k&=q_{\bm{\theta}_k}(\bm z_k(\bm x)).
  \label{eq:ket-sharing-mmoe}
\end{align}
$g_{k,m}$表示当前输入上专家$m$对任务$k$的混合权重，$\bm z_k$是任务表示，$q_{\bm{\theta}_k}$再将它读成输出。不同任务可以对同一组专家采用不同比例，同一任务也可以随输入调整比例，MMoE用这两种变化表达对共同映射的不同需求\scite{ket-ma2018-mmoe}。

训练时，任务损失沿这条路径更新自己的头和门控，也更新参与混合的专家；预测时重新计算输入对应的门控和专家输出。若改成所有任务共用一个仅依赖输入的门控，同一图像就得到同一混合表示；若改为Cross-stitch，系数则不再按输入重算。三者可以共用部件，但选择依据不同。

需要保留地区或任务差异时，还可在专家以外增加私有分支或校准参数。新增自由度也需要数据来估计，结构本身不保证个性化收益。样本多、更新频繁的任务可能先塑造共同专家，其他任务在有限预算内未必能充分调整；专家数量增加也可能只得到多份近似映射。

检验专家是否分工，应按任务及输入子群统计门控使用比例和集中程度，再观察删除或替换专家造成的预测变化，并与同预算固定共享比较。不同权重仅说明混合不同，不能直接给专家赋予语义。标准MMoE是稠密软混合，通常计算全部专家；要减少单次计算，还需稀疏选择和能跳过未选专家的执行机制，预算见第~\ref{sec:ket-sharing-budget}节。本节专家由任务共同训练，第~\ref{sec:ket-integration-conditional-modules}节则组合已经形成的模块。

\subsection{共享先验与参数关系的设计}
\label{sec:ket-sharing-prior}

病害模型数据少时，即使不复用类别模型的前向计算，也可能借助其参数估计来限制自己的自由度。这样共享的是参数可能取值的统计关系。先固定比较坐标：以下任务参数块$\bm w_k\in\mathbb R^d$必须同维、同坐标，可以是相同特征上的线性参数，或已对齐并按相同顺序展开的编码器块。异构的类别、病害和位置头仍各自保留，不将它们直接求平均；后面的式子省略这些私有部分。

最简单的联系是假设各参数块围绕一个共同中心变化。其他任务数据帮助确定中心，数据稀少的任务再在自身观测与这个中心之间取舍。

为把“围绕中心”写成概率模型，可采用下列简化层次高斯先验：
\begin{equation}
  \bm\mu\sim\mathcal N(\bm0,\kappa^2\bm I),
  \qquad
  \bm w_k\mid\bm\mu\sim
  \mathcal N(\bm\mu,\tau^2\bm I).
  \label{eq:ket-sharing-hierarchical-prior}
\end{equation}
其中$\bm I$是$d$维单位矩阵，$\kappa,\tau>0$为固定尺度：$\kappa$约束中心$\bm\mu$偏离零的幅度，$\tau$约束任务参数偏离中心的幅度。给定中心后，假设各任务参数条件独立，任务数据在给定各自参数后也条件独立。以$\mathcal L_k(\bm w_k)$记各任务负对数似然，最大化后验等价于把数据负对数似然与先验负对数密度相加后最小化：
\begin{equation}
  \sum_{k=1}^{K}\mathcal L_k(\bm w_k)
  +\frac{1}{2\tau^2}\sum_{k=1}^{K}
     \|\bm w_k-\bm\mu\|_2^2
  +\frac{1}{2\kappa^2}\|\bm\mu\|_2^2.
  \label{eq:ket-sharing-hierarchical-map}
\end{equation}
忽略固定尺度下的常数后，高斯先验留下两个二次惩罚：任务偏离中心要付出代价，中心远离零也要付出代价。对固定的各$\bm w_k$，令关于$\bm\mu$的导数为零，得到
\begin{equation}
  \bm\mu
  =\frac{\sum_{k=1}^{K}\bm w_k}
    {K+\tau^2/\kappa^2}.
\end{equation}
分母中$\tau^2/\kappa^2$使中心不是无约束的算术平均。中心固定以后，每项任务仍要拟合自己的数据；似然约束较弱的任务相对更受先验影响，参数被拉向中心，但不会被直接设为中心。这个共享机制成立与否取决于共同中心假设是否合理。

若任务分成多个群组，单一中心可能不合适。可以引入簇身份$c_k$，让任务参数围绕$\bm\mu_{c_k}$变化，并估计各簇的归属概率。Bakker和Heskes以贝叶斯多任务模型研究了共同先验、聚类及门控的联系\scite{ket-bakker2003-clustering}。上述单中心模型只解释收缩，不复现该论文完整推断。簇可依据可靠知识指定，也可由数据估计；后一种做法还需要足够样本区分簇，不能仅凭拟合灵活就认定分组可靠。

共同中心只约束参数接近，尚未描述任务之间哪些变化相伴发生。\term{任务关系学习}（Multi-Task Relationship Learning, MTRL）用任务参数间的协方差来表达这种联系。令$\bm W=[\bm w_1,\ldots,\bm w_K]\in\mathbb R^{d\times K}$，其每行对应同一参数坐标在$K$个任务中的取值。采用各行独立、每行服从零均值且协方差为$\bm\Omega\succ0$的高斯向量这一矩阵正态先验，负对数密度中与未知量有关的部分为
\begin{equation}
  \frac12\operatorname{tr}
    (\bm W\bm\Omega^{-1}\bm W^{\top})
  +\frac d2\log|\bm\Omega|.
  \label{eq:ket-sharing-matrix-prior}
\end{equation}
这里$\operatorname{tr}$为迹，$|\bm\Omega|$为行列式；第一项按协方差度量参数变化，第二项来自高斯归一化因子。当$\bm\Omega$未知时，第二项也会影响估计，不能视为常数删除。MTRL改用迹为1的约束固定关系矩阵尺度，并构造以下正则化准则：
\begin{align}
  \min_{\bm W,\bm\Omega}\quad&
  \sum_{k=1}^{K}\hat R_k(\bm w_k)
  +\frac{\lambda_1}{2}\|\bm W\|_{\mathrm F}^2
  \nonumber\\
  &+\frac{\lambda_2}{2}
    \operatorname{tr}(\bm W\bm\Omega^{-1}\bm W^{\top})
  \nonumber\\
  \text{满足}\quad&
  \bm\Omega\succeq0,
  \qquad \operatorname{tr}(\bm\Omega)=1.
  \label{eq:ket-sharing-mtrl}
\end{align}
其中$\lambda_1,\lambda_2>0$控制参数大小与关系惩罚，$\|\cdot\|_{\mathrm F}$为矩阵Frobenius范数。迹约束避免任意放大协方差来减小惩罚，半正定约束使其可作为协方差使用。这是改变估计准则后的凸正则化形式，不是删去行列式项就得到的高斯最大后验恒等式\scite{ket-zhang2010-mtrl}。

关系固定且任务损失凸时，可以求解带耦合正则的参数子问题；参数固定后，再求关系矩阵。令$\bm B=\bm W^{\top}\bm W$记录任务参数的内积，在$\bm W\ne\bm0$时，迹约束下的半正定最优关系为
\begin{equation}
  \bm\Omega^\star
  =\frac{\bm B^{1/2}}
    {\operatorname{tr}(\bm B^{1/2})}.
  \label{eq:ket-sharing-mtrl-omega}
\end{equation}
从特征分解可看出平方根的来源。将关系特征方向与$\bm B$对齐后，若$\bm B$的特征值为$b_j$，需在$\omega_j\geq0,\sum_j\omega_j=1$下最小化$\sum_jb_j/\omega_j$；正特征值方向上的拉格朗日条件给出$\omega_j\propto\sqrt{b_j}$。交替求参数和关系，使当前任务共同决定正则所用结构。训练结束后各任务仍用自己的参数预测，不必在部署中同时计算其他任务模型。

半正定允许零特征值，而式~\eqref{eq:ket-sharing-mtrl}写的是普通逆，边界含义必须补清。正定内部直接使用逆；边界使用伪逆，并要求$\bm W^{\top}$的值域包含在$\bm\Omega$的值域内，否则正则值定义为无穷大。数值上加入正定下界会改变允许集合，闭式更新也须调整。若$\bm W=\bm0$，所有迹为1的关系矩阵都有零关系惩罚，不能继续用式~\eqref{eq:ket-sharing-mtrl-omega}除以零。

协方差还允许相反变化，这与训练受损不同。对同特征坐标的两个任务，设$\rho$控制参数相关性，并采用
\begin{equation}
  \bm\Omega=\frac12
  \begin{bmatrix}1&\rho\\\rho&1\end{bmatrix},
  \qquad |\rho|<1.
\end{equation}
条件$|\rho|<1$使该矩阵正定且迹为1。将其逆代入关系惩罚，可以直接比较同向与反向参数的代价：
\begin{equation}
  \operatorname{tr}(\bm W\bm\Omega^{-1}\bm W^{\top})
  =
  \frac{2\bigl(\|\bm w_1\|_2^2+\|\bm w_2\|_2^2
  -2\rho\bm w_1^{\top}\bm w_2\bigr)}
  {1-\rho^2}.
  \label{eq:ket-sharing-negative-correlation}
\end{equation}
若$\bm w_2$接近$-\bm w_1$，在相同$|\rho|$下，负$\rho$比正$\rho$给予这组参数更小惩罚。两个任务若描述互补状态，一项参数估计可能帮助约束另一项的反向趋势。负相关描述这种统计联系；负迁移描述共同学习相对对照增加任务风险，不能因名称都含“负”便把两者等同。

参数关系依赖表示坐标。若一个任务翻转特征符号并同步改变其参数，或者重排隐藏单元而保持预测函数不变，表观参数相关性也可能变化；未对齐时便不能据此认定任务不同。协方差不是永久语义或因果关系，错误的先验形状、簇结构及不足的样本都可能产生不当收缩。第~\ref{sec:ket-transfer-prior}节先用历史任务准备先验再交给新任务，本节则由多个当前任务共同估计并使用，收益须在这些当前任务上分别检验。

\subsection{共享输出关系与约束的设计}
\label{sec:ket-sharing-outputs}

病害分支应观察哪个区域，可以由定位输出提供线索，这种联系比参数相关性更直接。不过，两个输出必须针对同一对象和同一裁剪坐标：多朵花的照片中，左侧类别不能约束右侧病害；叶片不可见时，没预测出叶斑也不能推出健康。在这些条件明确后，可以让输出参与另一任务的计算，或仅让输出共同满足训练约束。

让中间预测进入后续计算是一条共享路径。PAD-Net先预测深度、表面法向、轮廓和语义等中间任务，再将预测转换、汇合为特征以支持最终深度和语义输出；训练与测试均从图像生成这些中间量\scite{ket-xu2018-padnet}。

为说明花卉任务怎样传递信息，另设一个教学计算图，而不把它当作PAD-Net实验或默认接在MMoE之后。共同编码器$h$产生$\bm z$，$t_k$为接收任务形成基础表示，$u_j$给出中间预测，随后将其他任务的信息加入当前表示：
\begin{align}
  \bm z&=h(\bm x),
  \qquad \bm z_k=t_k(\bm z),
  \nonumber\\
  \widetilde y_j&=u_j(\bm z),
  \qquad
  \bm r_j=\phi_j(\widetilde y_j),
  \nonumber\\
  \bm v_k&=\bm z_k+
     \sum_{j\ne k}
     \bm a_{k,j}(\bm r_k,\bm r_j)
       \odot\psi_{k,j}(\bm r_j),
  \nonumber\\
  \hat y_k&=q_k(\bm v_k).
  \label{eq:ket-sharing-output-message}
\end{align}
$\phi_j$将中间预测$\widetilde y_j$编码成消息特征$\bm r_j$，$\psi_{k,j}$将发送特征对齐到接收任务。设接收网格的高、宽、通道数为$H_k,W_k,C_k$，则$\bm z_k,\bm v_k$及被相乘的两个消息张量都须为$\mathbb R^{H_k\times W_k\times C_k}$中的元素。图像级类别向量可先映射通道再广播，定位热图则须按同一裁剪坐标重采样；形状相同还不足以替代坐标对齐。

$\bm a_{k,j}$调节每个位置接收的消息，$\odot$为逐元素乘法，$q_k$把加入消息的$\bm v_k$转成最终输出。这样病害分支可利用定位热图与类别线索。训练以中间和最终监督共同约束这些可微路径；部署由输入图像重算中间预测，再完成依赖它的任务。

这条路径在使用时只能取得预测值，不能把训练时有的真实辅助标签当成免费输入。训练用真实区域、部署用预测区域，会改变下游面对的误差分布。概率图比硬标签多保留了输出分布信息，却不保证这些信息正确；应专门测试区域偏移、错误类别和低置信输入怎样影响下游。消息门控或注意力权重也不能未经校准便解释成置信概率。

如果不希望部署输出彼此依赖，可只在训练损失中惩罚不一致。为说明空间关系，暂将病害任务从图像级状态扩展为像素病变输出，并与同一对象区域对齐。病变概率$p_{\text{病}}(u\mid\bm x)$和区域概率$p_{\text{区}}(u\mid\bm x)$都定义在非空像素集合$\mathcal U$上；标签若规定病变只能位于该区域内，可以用区域外病变的平均软惩罚：
\begin{equation}
  L_{\mathrm{rel}}
  =\frac{1}{|\mathcal U|}
    \sum_{u\in\mathcal U}
    p_{\text{病}}(u\mid\bm x)
    \bigl[1-p_{\text{区}}(u\mid\bm x)\bigr],
  \label{eq:ket-sharing-output-constraint}
\end{equation}
该乘积在病变概率高、区域概率低的位置取大值，平均后避免像素数量直接决定惩罚尺度。在原监督损失上加$\lambda_{\mathrm{rel}}L_{\mathrm{rel}}$，其中$\lambda_{\mathrm{rel}}\geq0$，便可限制这种不一致。只预测图像级病害状态的模型不能直接使用它；而且病变全部为零或区域全部为一都可令惩罚为零，所以它不能替代标签监督。若区域定义错误或惩罚过重，还可能压制正确病变输出。

预测传递增加前向依赖与错误传播，一致性正则则只改变训练要求，部署可保持独立计算但不保证严格满足规则。第~\ref{sec:ket-integration-fixed-modules}节讨论已学好模块的串接，这里讨论关系怎样继续约束共同训练；规则取得与核验由知识获取主题承担。图~\ref{fig:ket-sharing-locations}用实线和虚线分别画计算与统计约束，读图时应据此判断哪些依赖在部署时仍需执行。

\begin{figure*}[htbp]
  \centering
  % Three sharing locations and a conditional-computation inset.
% Conceptual diagram only; all dimensions are in millimetres.
\begin{tikzpicture}[
  x=1mm,y=1mm,
  font=\figurefont\footnotesize,
  text=BookInk,
  >={Latex[length=1.8mm,width=1.3mm]},
  panel/.style={draw=BookRule,line width=0.6pt},
  model/.style={draw=BookTeal,fill=FigureModel!10,line width=0.8pt,
    minimum height=10mm,text width=11mm,align=center,inner sep=1mm},
  prior/.style={draw=BookBlue,fill=FigureInput!10,line width=0.8pt,
    minimum height=11mm,text width=27mm,align=center,inner sep=1mm},
  output/.style={text width=14mm,align=center,inner sep=0pt},
  relation/.style={draw=BookAmber,fill=FigureLoss!10,line width=0.8pt,
    minimum height=10mm,text width=40mm,align=center,inner sep=1mm},
  flow/.style={draw=BookInk,line width=0.8pt,-{Latex}},
  wire/.style={draw=BookInk,line width=0.8pt},
  stat/.style={draw=BookBlue,dashed,line width=0.8pt},
  constraint/.style={draw=BookAmber,dashed,line width=0.8pt},
  heading/.style={anchor=west,font=\figurefont\small}
]
  \path[use as bounding box] (0,5) rectangle (169,-160);
  \node[heading] at (1,0) {(a) 表示/模型参数 × 前向计算};
  \node[heading] at (59,0) {(b) 模型参数 × 训练先验};
  \node[heading] at (117,0) {(c) 预测输出 × 训练约束};
  \draw[panel] (0,-7) rectangle (53,-74);
  \draw[panel] (58,-7) rectangle (111,-74);
  \draw[panel] (116,-7) rectangle (169,-74);

  \node[model,text width=27mm] (shared) at (26.5,-20)
    {共同映射\\$h_{\bm{\theta}_{\mathrm s}}$};
  \node[model] (head1) at (8.5,-43) {$q_1$};
  \node[model] (head2) at (26.5,-43) {$q_2$};
  \node[model] (head3) at (44.5,-43) {$q_3$};
  \node[output] (outa1) at (8.5,-62) {类别\\$\hat y_1$};
  \node[output] (outa2) at (26.5,-62) {病害\\$\hat y_2$};
  \node[output] (outa3) at (44.5,-62) {定位\\$\hat y_3$};
  \draw[wire] (shared.south) -- (26.5,-31);
  \draw[flow] (26.5,-31) -| (head1.north);
  \draw[flow] (26.5,-31) -- (head2.north);
  \draw[flow] (26.5,-31) -| (head3.north);
  \draw[flow] (head1.south) -- (outa1.north);
  \draw[flow] (head2.south) -- (outa2.north);
  \draw[flow] (head3.south) -- (outa3.north);

  \node[prior] (prior) at (84.5,-20)
    {共同先验\\$\bm\mu,\bm\Omega$};
  \node[model] (modelb1) at (66.5,-43) {$f_{\bm w_1}$};
  \node[model] (modelb2) at (84.5,-43) {$f_{\bm w_2}$};
  \node[model] (modelb3) at (102.5,-43) {$f_{\bm w_3}$};
  \node[output] (outb1) at (66.5,-62) {类别\\$\hat y_1$};
  \node[output] (outb2) at (84.5,-62) {病害\\$\hat y_2$};
  \node[output] (outb3) at (102.5,-62) {定位\\$\hat y_3$};
  \draw[stat] (prior.south) -- (84.5,-31);
  \draw[stat] (84.5,-31) -| (modelb1.north);
  \draw[stat] (84.5,-31) -- (modelb2.north);
  \draw[stat] (84.5,-31) -| (modelb3.north);
  \draw[flow] (modelb1.south) -- (outb1.north);
  \draw[flow] (modelb2.south) -- (outb2.north);
  \draw[flow] (modelb3.south) -- (outb3.north);

  \node[model] (modelc1) at (124.5,-20) {$f_1$};
  \node[model] (modelc2) at (142.5,-20) {$f_2$};
  \node[model] (modelc3) at (160.5,-20) {$f_3$};
  \node[output] (outc1) at (124.5,-43) {类别\\$\hat y_1$};
  \node[output] (outc2) at (142.5,-43) {病害\\$\hat y_2$};
  \node[output] (outc3) at (160.5,-43) {定位\\$\hat y_3$};
  \node[relation] (relation) at (142.5,-63)
    {同对象预测输出约束\\$L_{\mathrm{rel}}$};
  \draw[flow] (modelc1.south) -- (outc1.north);
  \draw[flow] (modelc2.south) -- (outc2.north);
  \draw[flow] (modelc3.south) -- (outc3.north);
  \draw[constraint] (outc1.south) -- (124.5,-53) -- ($(relation.north)+(-10,0)$);
  \draw[constraint] (outc2.south) -- (relation.north);
  \draw[constraint] (outc3.south) -- (160.5,-53) -- ($(relation.north)+(10,0)$);

  \node[heading] at (0,-85) {(d) 表示/模型参数 × 条件前向：专家由各任务共同训练};
  \node[prior,text width=16mm] (input) at (12,-116)
    {同一图像\\$\bm x$};
  \node[model,text width=25mm] (expert1) at (54,-101)
    {专家1：$\bm e_1(\bm x)$};
  \node[model,text width=25mm] (expert2) at (54,-116)
    {专家2：$\bm e_2(\bm x)$};
  \node[model,text width=25mm] (expert3) at (54,-131)
    {专家3：$\bm e_3(\bm x)$};
  \node[prior,text width=28mm] (gate) at (54,-151)
    {任务$k$的门控\\$\bm g_k(\bm x)$};
  \node[model,text width=24mm,minimum height=16mm] (mix) at (111,-116)
    {按权重混合\\$\bm z_k$};
  \node[model,text width=24mm,minimum height=16mm] (task) at (151,-116)
    {任务头$q_k$\\预测输出$\hat y_k$};
  \node[output,text width=30mm] at (151,-134)
    {$k=1,2,3$\\类别、病害、定位};
  \draw[wire] (input.east) -- (28,-116);
  \draw[flow] (28,-116) |- (expert1.west);
  \draw[flow] (28,-116) -- (expert2.west);
  \draw[flow] (28,-116) |- (expert3.west);
  \draw[flow] (28,-116) |- (gate.west);
  \draw[flow] (expert1.east) -- (mix.north west);
  \draw[flow] (expert2.east) -- (mix.west);
  \draw[flow] (expert3.east) -- (mix.south west);
  \draw[flow] (gate.east) -- (111,-151) -- (mix.south);
  \node[anchor=west] at (114,-145) {混合权重};
  \draw[flow] (mix.east) -- (task.west);
\end{tikzpicture}

  \caption{共同计算、参数先验与输出约束的共享位置。上排固定类别、病害和定位任务，省略重复输入；实线箭头是前向计算，虚线是先验联系或训练约束。先验面板只展开同维同坐标参数块$\bm w_k$的联系，未画异构输出头。输出面板示意训练约束，并未画出预测传递所需的额外前向路径。下排将共同计算展开为条件混合：同一图像送入专家与任务$k$的门控，再由任务权重混合专家输出；此处$\bm g_k(\bm x)$是由$g_{k,m}$组成的门控向量，不是共享梯度。三个任务分别设置门控和头，损失反传到对应部件及共同专家；反向路径未画。稠密门控仍需计算全部专家。}
  \label{fig:ket-sharing-locations}
\end{figure*}
\FloatBarrier

\section{联合学习目标的协调}
\label{sec:ket-sharing-objectives}

三个任务使用同一骨干后，损失如何合成就决定它的一次更新。这里暂时固定任务分组、结构与采样，分别看损失尺度、梯度大小和梯度方向。它们不能互相代替：数值大的损失未必变化快，范数大的梯度未必对应重要任务，方向冲突也只说明当前更新的局部问题，尚不能认定长期负迁移。

可以先给任务损失加权，再求一个标量目标的梯度；也可以保留每项任务梯度，直接求联合方向。后一种求法也可能得到权重，因而两种操作不是互斥类别。若同一次求导中系数相同且作为常量，加权损失与加权梯度相等；需要比较的是系数如何确定、何时重估，以及所得更新满足什么条件。

\subsection{任务权重的确定}
\label{sec:ket-sharing-weights}

假设当前批次包含所有任务监督，先以非负权重$w_k$表示每项损失对更新的贡献，得到加权联合目标：
\begin{equation}
  L_{\mathrm{joint}}
  =\sum_{k=1}^{K}w_kL_k,
  \qquad w_k\geq0.
  \label{eq:ket-sharing-weighted-loss}
\end{equation}
权重固定时，共享梯度为$\sum_k w_k\bm g_k$。放大某项贡献可以表达优先级，也可以补偿数值尺度，二者应分别记录。例如，对同一物理位置预测，把坐标由米换成厘米再计算平方误差，损失及其参数梯度都扩大$10^4$倍；类别交叉熵不变，定位责任也没有改变。因此，要先明确单位、归一化方式及按样本还是像素平均，再讨论收益取舍。

权重还应对应前面的交付责任。共同部署时要保护各项最低标准，只部署类别时以类别验证结果选辅助权重，地区副本则需说明等权还是按服务量加权。固定权重给出明确可复现的目标，却未必在所有训练阶段都形成合适作用；动态调整能根据训练状态变化，但不能自行决定业务优先级。

评价调权方法时，应与同数据、同结构、同基础优化器的简单联合训练比较。在预定搜索预算内调整固定权重或采样比例，并保持停止规则一致，才能判断专用协调器增加了什么。单任务对照还改变了共享结构，不能单独隔离调权作用。权重搜索、逐任务求梯度和额外子问题求解都计入费用；采样本身改变目标的方式见第~\ref{sec:ket-sharing-sampling}节。

若希望下降较慢的任务得到更多作用，可以读取近期损失轨迹。\term{动态权重平均}（Dynamic Weight Average, DWA）将任务损失的相邻阶段比值转成权重。这里$L_k^{(t)}$采用阶段内平均损失，而非单个随机批次值；分母需为正。令
\[
  r_k^{(t-1)}=\frac{L_k^{(t-1)}}{L_k^{(t-2)}},
  \qquad
  w_k^{(t)}
  =K\frac{\exp(r_k^{(t-1)}/T)}
          {\sum_j\exp(r_j^{(t-1)}/T)},
\]
$T>0$为平滑权重差异的温度，权重和为$K$；尚无两个历史阶段时使用预定初值。损失下降慢时比值$r_k$较大，权重也较大，但尺度突变、批次噪声或接近零的分母都可能使这个依据失真。DWA处理的是相对下降速度，不判断任务重要性，也不消除梯度方向冲突\scite{ket-liu2019-mtan}。

另一种依据是解释观测误差的噪声尺度。\term{不确定性加权}（uncertainty weighting）把任务权重与所假设的观测模型联系起来。暂取$d_k$维连续输出，令任务预测$f_k(\bm x)$与目标之间的残差为独立高斯噪声：
\begin{equation}
  y_k=f_k(\bm x)+\bm\varepsilon_k,
  \qquad
  \bm\varepsilon_k\sim
  \mathcal N(\bm0,\sigma_k^2\bm I).
\end{equation}
这里$\sigma_k>0$对同一任务的全部输入保持相同，称为任务级同方差噪声尺度，并不是每张图的置信分数。把逐样本平方残差范数的一半在批次上平均，记为$L_k^{\mathrm{reg}}$；再以$s_k=\log\sigma_k^2$参数化正方差。高斯负对数密度展开后，平均负对数似然为
\begin{equation}
  \mathcal L_k^{\mathrm{reg}}
  =e^{-s_k}L_k^{\mathrm{reg}}
   +\frac{d_k}{2}s_k+C.
  \label{eq:ket-sharing-uncertainty-regression}
\end{equation}
$C$收集与模型参数和$s_k$无关的常数。残差项按逆方差加权，方差越大作用越小；对数方差项则阻止模型仅靠增大方差来忽略任务。$d_k/2$来自$d_k$个输出的高斯归一化因子，若改成按维度平均，应相应缩放整项似然。各任务分别取样本平均也属于任务归一化选择，样本数不同时不等于直接对全部观测似然求和。

分类的概率归一化方式不同，不能照搬回归式。记$\bm f_k$为未归一化类别分数，$f_{k,c}$为类别$c$的分量，$L_k^{\mathrm{CE}}$为原分数的交叉熵。把概率改为$\operatorname{softmax}(e^{-s_k}\bm f_k)$后，展开单样本负对数概率，可得
\begin{align}
  \mathcal L_k^{\mathrm{cls}}
  ={}&e^{-s_k}L_k^{\mathrm{CE}}
  +\log\sum_c\exp(e^{-s_k}f_{k,c})
  \nonumber\\
  &-e^{-s_k}\log\sum_c\exp(f_{k,c}),
  \label{eq:ket-sharing-uncertainty-classification}
\end{align}
后两个对数项来自softmax归一化，不能任意消去。Kendall等近似处理这一归一化因子，得到常用的$e^{-s_k}L_k^{\mathrm{CE}}+\tfrac12s_k$；这个简化不是任意类别分数下的恒等式\scite{kendall2018}。使用时应说明采用的是精确分类概率还是该近似目标。

选定似然形式后，可从$s_k=0$开始，逐批同时更新模型与噪声参数。为看清噪声更新的含义，将其中一项写成$e^{-s_k}L_k+c_ks_k$，$c_k>0$为相应对数项系数；固定模型参数求导，有
\begin{equation}
  \frac{\partial\mathcal L_k}{\partial s_k}
  =-e^{-s_k}L_k+c_k.
\end{equation}
在$L_k>0$时，令导数为零得到$\sigma_k^2=L_k/c_k$，因此权重对应当前误差尺度，并不独立表达优先级。误差趋零时方差可能走向零边界，需另加数值下界或先验。还应检查噪声假设是否适当、困难任务是否过早被降权，以及任务平均是否掩盖稀缺群体误差。

损失归一化以后，共享梯度范数仍可能相差很大。梯度归一化（Gradient Normalization, GradNorm）根据梯度大小与相对学习进度调权，使下降较慢的任务获得较大的目标梯度作用。先选择测量块$\bm{\theta}_{\mathrm b}$，例如任务分支前的共享层，并以$\bm g_{k,\mathrm b}^{(t)}$表示任务在该块的未加权梯度。将当前权重视为正的常量，测得
\begin{equation}
  G_k^{(t)}
  =\left\|\nabla_{\bm{\theta}_{\mathrm b}}
     \bigl(w_k^{(t)}L_k^{(t)}\bigr)\right\|_2
  =w_k^{(t)}\|\bm g_{k,\mathrm b}^{(t)}\|_2,
  \label{eq:ket-sharing-gradnorm-size}
\end{equation}
保持权重为正可令$w_k=\exp(v_k)$，其中$v_k$为无约束的调权变量，也可在更新后投影到正数域；不应在无约束更新后默认正性仍成立。为了区分绝对损失与下降进度，先用非零初始损失$L_k^{(0)}$归一化，再与所有任务的平均进度相比：
\begin{equation}
  q_k^{(t)}=\frac{L_k^{(t)}}{L_k^{(0)}},
  \qquad
  r_k^{(t)}
  =\frac{q_k^{(t)}}{K^{-1}\sum_jq_j^{(t)}}.
  \label{eq:ket-sharing-gradnorm-rate}
\end{equation}
$q_k$为当前损失占自身初值的比例，$r_k>1$表示这一比例高于任务平均，即下降较慢，而非绝对损失较大。初值为零或极小时应预先另设有效参照；若所有当前损失都为零，平均进度也为零，本步应停止这项调权而非继续相除。

记平均加权范数为$\bar G=K^{-1}\sum_kG_k$。以$\bar G r_k^\alpha$作为任务$k$的目标范数，$\alpha\geq0$决定对慢任务倾斜多少；让测量范数接近这个目标，可最小化
\begin{equation}
  L_{\mathrm{grad}}
  =\sum_{k=1}^{K}
  \left|G_k-
    \operatorname{stopgrad}(\bar G r_k^\alpha)\right|.
  \label{eq:ket-sharing-gradnorm-objective}
\end{equation}
绝对值衡量测量值与目标值的距离，$\operatorname{stopgrad}$要求本次调权把目标当作常量。已取得的任务梯度也固定，只对$w_k$求导；否则目标随权重一起变化，或直接用这个目标训练表示，就改变了原有职责。模型仍按加权任务损失更新，权重更新后保持正值并归一化为$w_k\leftarrow Kw_k/\sum_jw_j$。这使GradNorm同时利用梯度大小和相对学习速度\scite{ket-chen2018-gradnorm}。

假设教学进度$q_k=(0.2,0.5,0.8)$，其均值为$0.5$，故$r_k=(0.4,1,1.6)$。再取$\bar G=2,\alpha=1$，三个目标范数为$(0.8,2,3.2)$，下降慢的定位获得较大目标。这是在调整训练作用，不证明定位应享有更高优先级。若下降慢是标签噪声造成，持续追赶可能浪费计算；测量层、初值及批次波动也影响权重。噪声加权与GradNorm依据不同，都须用逐任务验证检查结果。

\subsection{联合更新方向的构造}
\label{sec:ket-sharing-directions}

调成相近的梯度范数以后，两项任务仍可能要求相反变化。下面固定私有参数并省略阶段上标，以$\bm d$表示共享梯度组合，实际参数更新为$\bm{\theta}_{\mathrm s}^{+}=\bm{\theta}_{\mathrm s}-\eta\bm d$，$\eta>0$。要判断这一步帮助哪项任务，可沿该方向展开局部Lipschitz梯度的任务损失：
\begin{equation}
  L_k(\bm{\theta}_{\mathrm s}-\eta\bm d)
  =L_k(\bm{\theta}_{\mathrm s})
   -\eta\bm g_k^{\top}\bm d+O(\eta^2).
  \label{eq:ket-sharing-direction-effect}
\end{equation}
一阶项是$-\eta\bm g_k^{\top}\bm d$，内积为正时足够小的一步会降低当前损失。由这个判据，可以移除冲突分量，也可以直接找全部任务的共同下降方向，还可以在共同下降方向中选择收益取舍。这三种要求可叠加，并非互斥且穷尽的分类。

\subsubsection{基于冲突消解的方向构造}
\label{sec:ket-sharing-pcgrad}

若合成前先删去某项梯度中妨碍另一任务下降的分量，便得到投影冲突梯度（projecting conflicting gradients, PCGrad）的处理思路，即把冲突副本投到另一梯度的正交平面。为任务$k$复制原始共享梯度，记为$\widetilde{\bm g}_k=\bm g_k$；遍历其他任务$j$时，用当前副本与原始$\bm g_j$的内积判断冲突。内积为负且$\bm g_j$非零时，作投影：
\begin{equation}
  \widetilde{\bm g}_k\leftarrow
  \widetilde{\bm g}_k-
  \frac{\widetilde{\bm g}_k^{\top}\bm g_j}
       {\|\bm g_j\|_2^2}\bm g_j.
  \label{eq:ket-sharing-pcgrad-projection}
\end{equation}
所减项是当前副本在$\bm g_j$方向上的分量，删去后与$\bm g_j$的内积为零。无冲突时不改，零梯度跳过，以免除零。每项副本按随机顺序检查其他任务，完成后求和或按约定取平均；私有参数仍用本任务监督更新。投影参照始终是其他任务的原始梯度，不是对方已经修改的副本\scite{ket-yu2020-pcgrad}。

用同一组二维教学损失，可以复算投影、共同下降与协商三种构造。令$\bm\theta=(\theta_1,\theta_2)^{\top}$为两个共享自由度，两项任务各希望接近不同中心：
\begin{align}
  L_1(\bm\theta)
    &=\tfrac12\|\bm\theta-(-1,0)^{\top}\|_2^2,
  \nonumber\\
  L_2(\bm\theta)
    &=\tfrac12\|\bm\theta-(2,-2)^{\top}\|_2^2.
  \label{eq:ket-sharing-toy-losses}
\end{align}
这些二次损失只模拟共享块竞争，不是花卉测量。零点梯度分别为$\bm g_1=(1,0)^{\top}$、$\bm g_2=(-2,2)^{\top}$，内积$-2$触发冲突处理。直接求和为$\bm d_{\mathrm{sum}}=(-1,2)^{\top}$，与$\bm g_1$内积为$-1$，因而小步更新会提高任务1损失，即使它可能降低总损失。

将这两个梯度代入式~\eqref{eq:ket-sharing-pcgrad-projection}，各自删去相对另一任务的冲突分量：
\begin{align}
  \widetilde{\bm g}_1
    &=(1,0)^{\top}-\frac{-2}{8}(-2,2)^{\top}
      =(0.5,0.5)^{\top},
  \nonumber\\
  \widetilde{\bm g}_2
    &=(-2,2)^{\top}-\frac{-2}{1}(1,0)^{\top}
      =(0,2)^{\top}.
\end{align}
求和后$\bm d_{\mathrm{PC}}=(0.5,2.5)^{\top}$，与原梯度内积为$0.5,4$，两项一阶变化均向下。再固定步长$\eta=0.1$，将实际更新参数代回二次损失，而不只检查一阶项，得到
\begin{align}
  (L_1,L_2)(\bm0)&=(0.5,4),
  \nonumber\\
  (L_1,L_2)(-\eta\bm d_{\mathrm{sum}})
    &=(0.625,3.425),
  \nonumber\\
  (L_1,L_2)(-\eta\bm d_{\mathrm{PC}})
    &=(0.4825,3.6325).
  \label{eq:ket-sharing-pcgrad-numbers}
\end{align}
直接求和降低了总损失，却把任务1从$0.5$增为$0.625$；投影后任务1为$0.4825$，两项均下降。这只比较当前损失和步长，两方向长度不同，不能由此排出一般预测性能次序。

多任务逐对投影还存在顺序影响：后一次投影可能重新改变前面已处理的角度，不能把结果当成同时满足全部约束的精确投影。批次负内积也可能含估计噪声。分别求梯度和追加内积增加训练费用，所得改善仍是当前更新的行为，不保证最终所有任务超过同预算单任务对照。

\subsubsection{基于共同下降的方向构造}
\label{sec:ket-sharing-mgda}

若不只处理一对冲突，而要求本步同时降低全部当前任务损失，可以在所有任务梯度的凸组合中寻找方向。\term{多梯度下降算法}（multiple gradient descent algorithm, MGDA）采用非负且和为1的系数，以梯度凸包为候选集，选择距零点最近的组合：
\begin{align}
  \bm\alpha^\star
  &=\argmin_{\substack{\alpha_k\geq0\\\sum_k\alpha_k=1}}
  \frac12\left\|\sum_{k=1}^{K}\alpha_k\bm g_k\right\|_2^2,
  \nonumber\\
  \bm d^\star&=\sum_k\alpha_k^\star\bm g_k.
  \label{eq:ket-sharing-mgda}
\end{align}
平方范数把“距原点最近”写成可求解的凸二次子问题。$\alpha_k^\star$是这个子问题的系数，不是预先指定的业务权重。Sener与Koltun将该构造用于共享网络的多目标更新，并研究降低其求解费用的近似\scite{ket-sener2018-mgda}。

最短组合为何能保证一阶共同下降？从$\bm d^\star$朝凸包任一顶点$\bm g_k$移动仍在候选集内，最优性要求这个移动的单侧方向导数非负。对平方范数求导便有
\begin{equation}
  (\bm g_k-\bm d^\star)^{\top}\bm d^\star\geq0,
  \qquad
  \bm g_k^{\top}\bm d^\star
  \geq\|\bm d^\star\|_2^2.
  \label{eq:ket-sharing-common-descent}
\end{equation}
非零$\bm d^\star$使右侧严格为正，所以所有任务沿$-\bm d^\star$都一阶下降。要把它换成实际有限步下降，还需控制二阶余项；若各梯度在所走邻域内有光滑常数$\beta_k>0$，下降引理给出
\begin{equation}
  L_k(\bm\theta-\eta\bm d)
  \leq L_k(\bm\theta)
  -\eta\bm g_k^{\top}\bm d
  +\frac{\beta_k\eta^2}{2}\|\bm d\|_2^2.
\end{equation}
其中$\bm\theta$表示当前共享块。对$\bm d=\bm d^\star\ne0$，取$0<\eta<2/\max_k\beta_k$且保持在该邻域内，二次上界便小于原损失。任意大步长不继承这项结论；光滑与凸优化基础见\BookChapterRef{chap:convex-learning-and-stability}{卷3《基础、理论与可信性》第7章“一般学习的可学习性”}。

回到式~\eqref{eq:ket-sharing-toy-losses}的零点，以$\alpha_1=a,\alpha_2=1-a$参数化线段，组合为$(3a-2,2-2a)^{\top}$。在$0\leq a\leq1$上最小化平方范数，导数为零得到$a=10/13$，从而
\begin{equation}
  \bm d^\star=(4/13,6/13)^{\top},
  \qquad
  \bm g_1^{\top}\bm d^\star
  =\bm g_2^{\top}\bm d^\star
  =4/13.
  \label{eq:ket-sharing-mgda-numbers}
\end{equation}
两个内积都为$4/13>0$，说明当前方向满足共同下降判据。但移到两中心中点$\bm\theta=(0.5,-1)^{\top}$后，梯度为$(1.5,-1)^{\top}$及其相反向量，等权组合为零。任何方向与它们的内积互为相反数，故不能同时严格为正。

一般无约束且可微的情形下，零属于梯度凸包便称为\term{帕累托驻点}（Pareto stationary point），即不存在使全部目标严格下降的一阶方向。\term{帕累托有效}（Pareto efficient）则比较实际可行解：不存在另一解使全部目标不差且至少一项更好。前者是局部一阶条件，不保证后者；非凸网络可能停在较差驻点。上述凸二次中点恰为有效取舍，但仅计算共享块时，判断仍限于固定私有参数的该块，任务头还可更新。这两项性质也都不等于超越单任务基线。

求精确组合昂贵时，可在系数单纯形上用Frank--Wolfe迭代。当前$\bm d=\sum_k\alpha_k\bm g_k$，先按线性化目标选择内积$\bm g_j^{\top}\bm d$最小的顶点$\bm e_j$，再在$0\leq\gamma\leq1$上搜索并更新$\bm\alpha\leftarrow(1-\gamma)\bm\alpha+\gamma\bm e_j$。用当前范数与最小顶点内积之差衡量求解误差：
\begin{equation}
  \Delta=\|\bm d\|_2^2-\min_k\bm g_k^{\top}\bm d.
\end{equation}
若$\Delta<\|\bm d\|_2^2$，则最小顶点内积也为正，当前近似方向仍共同一阶下降。只设固定迭代次数而不检查误差，不能沿用精确解的判据。

另一种节省方式是减少所保存梯度的维数。MGDA-UB在共享表示上求一个上界问题：设$\bm z=h_{\bm{\theta}_{\mathrm s}}(\bm x)$，$\bm r_k=\nabla_{\bm z}L_k$，$\bm J=\partial\bm z/\partial\bm{\theta}_{\mathrm s}$为雅可比，则链式法则给出$\bm g_k=\bm J^{\top}\bm r_k$。以相同凸系数构造参数方向时，真实下降判据中的内积为
\begin{equation}
  \bm g_k^{\top}\bm d
  =\bm r_k^{\top}\bm J\bm J^{\top}
    \sum_j\alpha_j\bm r_j.
  \label{eq:ket-sharing-mgda-feature-metric}
\end{equation}
因此，直接最小化表示梯度范数并不是同一个度量。其计算依据是谱范数对任意$\bm r(\bm\alpha)=\sum_k\alpha_k\bm r_k$给出的上界：
\[
  \left\|\sum_k\alpha_k\bm g_k\right\|_2
  =\|\bm J^\top\bm r(\bm\alpha)\|_2
  \leq\|\bm J\|_2\,\|\bm r(\bm\alpha)\|_2.
\]
$\|\bm J\|_2$在同一参数状态和批次下不随组合系数改变，故可通过表示范数控制参数范数的上界。但上界最小不保证真实范数最小，更不直接保证各任务真实内积正。若$\bm J\bm J^{\top}$在这些梯度张成的子空间上等于$c\bm I$且$c>0$，两种度量只差正常数，可继承下降关系；一般情况须检查方法假设或复查真实内积。即使真实批次内积全正，结论也只针对当前批次，不承诺总体风险逐步降低。

\subsubsection{基于协商收益的方向构造}
\label{sec:ket-sharing-nash}

共同下降方向往往不止一个，还需决定怎样分配一阶收益。Nash-MTL以不更新为参考状态，让任务按预定效用进行协商\scite{ket-navon2022-nash}。为比较方向而排除任意放大，固定$\|\bm d\|_2\leq1$，省略共同步长因子，以$u_k(\bm d)=\bm g_k^{\top}\bm d$为任务收益。只接受各项收益为正的方向，再最大化对数收益和：
\begin{equation}
  \max_{\substack{\|\bm d\|_2\leq1\\
                  \bm g_k^{\top}\bm d>0\ \forall k}}
  \sum_{k=1}^{K}\log(\bm g_k^{\top}\bm d).
  \label{eq:ket-sharing-nash-objective}
\end{equation}
正效用要求每项相对零更新都一阶受益；对数和等价于效用乘积，不能用某项负收益交换其他项正收益。这个准则比较的是给定损失的局部线性收益，不是直接比较最终验证分数。若不存在严格共同下降方向，可行集为空，应停止套用该协商目标。

把共享梯度列成$\bm G=[\bm g_1,\ldots,\bm g_K]$，其Gram矩阵$\bm H=\bm G^{\top}\bm G$记录任务两两内积。存在严格可行方向时，向外缩放增加全部效用，所以最优解位于单位球面。用$\lambda$作球约束乘子，关于方向求导有
\begin{equation}
  \sum_k\frac{\bm g_k}{\bm g_k^{\top}\bm d}
  =\lambda\bm d.
\end{equation}
与单位方向$\bm d$取内积后，左侧各项都是1，得到$\lambda=K$。再定义正系数$a_k=1/(\sqrt K\,\bm g_k^{\top}\bm d)$，驻点关系写成$\bm d=\bm G\bm a/\sqrt K$。将它代回效用，得到系数必须满足的方程：
\begin{equation}
  \bm H\bm a=\bm a^{-1},
  \qquad
  (\bm a^{-1})_k=1/a_k.
  \label{eq:ket-sharing-nash-system}
\end{equation}
左乘$\bm a^{\top}$即得$\|\bm G\bm a\|_2^2=K$，因此$\bm d_{\mathrm N}=\bm G\bm a/\sqrt K$为单位方向，任务效用为$1/(\sqrt K a_k)$。这些系数共同满足协商关系，不是任意加权。梯度线性独立使$\bm H$正定，是下面求法的充分条件；梯度秩亏、近乎相关或没有严格共同下降方向时，不能不检查条件便求逆。

原论文使用凸代理与凸凹过程近似求权重。这里为手算另给一个教学求法，不声称复现其实现。在$\bm H\succ0$、$\bm a>0$下，考虑二次项与对数屏障的和，其梯度为零正好是式~\eqref{eq:ket-sharing-nash-system}：
\begin{equation}
  \psi(\bm a)
  =\tfrac12\bm a^{\top}\bm H\bm a
   -\sum_k\log a_k,
  \qquad a_k>0.
  \label{eq:ket-sharing-nash-newton}
\end{equation}
该函数严格凸，梯度$\bm r=\bm H\bm a-\bm a^{-1}$，海森矩阵$\bm H+\operatorname{diag}(a_k^{-2})$正定。以全正向量初始化，求解牛顿方向的线性方程：
\[
  \bigl[\bm H+\operatorname{diag}(a_k^{-2})\bigr]
  \bm\delta=-\bm r
\]
所得$\bm\delta$是系数增量。通过回溯选择步长，使全部$a_k$继续为正且$\psi$减小；以方程残差达到预设容差为停止条件，再归一化$\bm G\bm a$。求解误差存在时，这个方向也只是协商解的近似。

重估系数需要取得逐任务梯度、形成Gram矩阵并求解$K$维子问题。复用多步前的系数可省重估费用，但不再精确满足当前梯度的协商方程。

沿用零点教学梯度，$\bm H=\left[\begin{smallmatrix}1&-2\\-2&8\end{smallmatrix}\right]$。两行方程分别乘$a_1,a_2$，得$a_1(a_1-2a_2)=1$和$a_2(-2a_1+8a_2)=1$；相减后$a_1^2=8a_2^2$。选择全正根便有
\begin{align}
  a_2&=(8-4\sqrt2)^{-1/2},
  \qquad a_1=2\sqrt2\,a_2,
  \nonumber\\
  \bm d_{\mathrm N}&\approx(0.3827,0.9239)^{\top}.
\end{align}
前面MGDA方向长度不是1，不能直接与单位Nash方向比较收益。将其归一化为$\bm d_{\mathrm M}=(2,3)^{\top}/\sqrt{13}$，再对同一梯度计算两种单位方向的效用：
\begin{align}
  (u_1,u_2)(\bm d_{\mathrm M})
    &\approx(0.5547,0.5547),
  \nonumber\\
  (u_1,u_2)(\bm d_{\mathrm N})
    &\approx(0.3827,1.0824).
  \label{eq:ket-sharing-nash-numbers}
\end{align}
在这个构造中，归一化MGDA给两项相同效用，Nash给任务1较小的正效用、给任务2较大效用。后者乘积从$4/13\approx0.3077$增为$\sqrt2-1\approx0.4142$，满足协商选择，却没有逐项超过MGDA。不同准则选择不同取舍，并不形成每项收益一致的排名。

图~\ref{fig:ket-sharing-gradient-geometry}固定同一坐标与原梯度，分别画投影后的组合、凸包最近点和协商单位方向。观察箭头时应区分梯度组合$\bm d$与参数移动：优化沿$-\eta\bm d$执行，不能把图中箭头本身当作实际位移。

\begin{figure*}[htbp]
  \centering
  % Teaching geometry from eq:ket-sharing-toy-losses; no experimental data.
% Source gradients are g1=(1,0), g2=(-2,2), not (-2,3).
% All panels use 11.5 mm per coordinate unit; canvas is 169 x 74 mm.
% Arrows show gradients or their combinations d; displacement is -eta*d.
\begin{tikzpicture}[
  x=1mm,y=1mm,
  font=\figurefont\footnotesize,
  text=BookInk,
  >={Latex[length=1.7mm,width=1.2mm]},
  axis/.style={draw=BookMuted,line width=0.6pt,-{Latex}},
  tick/.style={draw=BookMuted,line width=0.5pt},
  original/.style={draw=FigureInput,line width=0.9pt,-{Latex}},
  projected/.style={draw=FigureModel,line width=1pt,-{Latex}},
  result/.style={draw=FigureLoss,line width=1.1pt,-{Latex}},
  guide/.style={draw=BookMuted,dashed,line width=0.7pt},
  leader/.style={draw=BookMuted,line width=0.5pt},
  heading/.style={anchor=west,font=\figurefont\footnotesize\bfseries},
  label/.style={inner sep=0.6mm},
  note/.style={align=center,inner sep=0pt}
]
  \path[use as bounding box] (0,0) rectangle (169,74);
  \node[heading] at (0,70) {(a) PCGrad：投影};
  \node[heading] at (58,70) {(b) MGDA：凸包最短点};
  \node[heading] at (116,70) {(c) 单位方向与共同下降};

  % Identical origins, ranges and physical scales in all three panels.
  \foreach \offset in {0,58,116} {
    \begin{scope}[shift={(\offset,0)}]
      \begin{scope}[shift={(29,25)},x=11.5mm,y=11.5mm]
        \draw[axis] (-2.25,0) -- (1.45,0);
        \draw[axis] (0,-1.2) -- (0,2.9);
        \foreach \v in {-2,-1,1} {
          \draw[tick] (\v,-0.045) -- (\v,0.045);
          \node[label,anchor=north] at (\v,-0.09) {$\v$};
        }
        \foreach \v in {-1,1,2} {
          \draw[tick] (-0.045,\v) -- (0.045,\v);
          \node[label,anchor=east] at (-0.09,\v) {$\v$};
        }
        \node[label,anchor=north east] at (-0.06,-0.06) {$0$};
      \end{scope}
      \node[label,anchor=west] at (46.5,25) {$d_1$};
      \node[label,anchor=south] at (29,60) {$d_2$};
    \end{scope}
  }

  \begin{scope}[shift={(29,25)},x=11.5mm,y=11.5mm]
    \coordinate (p1) at (0.5,0.5);
    \coordinate (p2) at (0,2);
    \draw[guide] (1,0) -- (p1);
    \draw[guide] (-2,2) -- (p2);
    \draw[original] (0,0) -- (1,0);
    \draw[original] (0,0) -- (-2,2);
    \draw[projected] (0,0) -- (p1);
    \draw[projected] (0,0) -- (p2);
    \draw[result] (0,0) -- (0.5,2.5);
    % The two removed components are orthogonal to the projected vectors.
    \draw[leader] (0.42,0.42) -- (0.50,0.34) -- (0.58,0.42);
    \draw[leader] (-0.14,2) -- (-0.14,1.86) -- (0,1.86);
  \end{scope}
  \node[label,anchor=west] at (42,29) {$\bm g_1$};
  \node[label,anchor=south] at (6,49) {$\bm g_2$};
  \node[label,anchor=west] at (38,35) {$\widetilde{\bm g}_1$};
  \draw[leader] (p1) -- (38,34);
  \node[label,anchor=south] at (21,54) {$\widetilde{\bm g}_2$};
  \draw[leader] (p2) -- (21,54);
  \node[label,anchor=west] at (37,56) {$\bm d_{\mathrm{PC}}$};
  \node[note] at (26.5,4) {$\bm d_{\mathrm{PC}}=(0.5,2.5)^{\top}$};

  \begin{scope}[shift={(87,25)},x=11.5mm,y=11.5mm]
    \coordinate (mgda) at ({4/13},{6/13});
    \draw[draw=BookInk,line width=1pt] (-2,2) -- (1,0);
    \fill[FigureInput] (-2,2) circle[radius=0.9mm];
    \fill[FigureInput] (1,0) circle[radius=0.9mm];
    \draw[projected] (0,0) -- (mgda);
    \fill[FigureModel] (mgda) circle[radius=0.7mm];
    % Perpendicularity to g1-g2=(3,-2) identifies the closest point.
    \draw[leader]
      ({4/13+0.16*3/sqrt(13)},{6/13-0.16*2/sqrt(13)})
      -- ({4/13+0.16/sqrt(13)},{6/13-0.16*5/sqrt(13)})
      -- ({4/13-0.16*2/sqrt(13)},{6/13-0.16*3/sqrt(13)});
  \end{scope}
  \node[label,anchor=south] at (64,49) {$\bm g_2$};
  \node[label,anchor=west] at (99,29) {$\bm g_1$};
  \node[label,anchor=west] at (96,37) {$\bm d^\star$};
  \draw[leader] (mgda) -- (96,36);
  \node[note] at (84.5,4) {$\bm d^\star=(4/13,6/13)^{\top}$};

  \begin{scope}[shift={(145,25)},x=11.5mm,y=11.5mm]
    % g1.d>0 and g2.d>0 give d1>0 and d2>d1: angles 45 to 90 deg.
    \fill[FigureModel!12] (0,0) -- (45:1)
      arc[start angle=45,end angle=90,radius=11.5mm] -- cycle;
    \draw[draw=BookMuted,dotted,line width=0.7pt] (0,0) circle[radius=11.5mm];
    \draw[guide] (45:1) -- (0,0) -- (90:1);
    \coordinate (unitm) at ({2/sqrt(13)},{3/sqrt(13)});
    \coordinate (unitn) at ({sin(22.5)},{cos(22.5)});
    \draw[projected] (0,0) -- (unitm);
    \draw[result,dashed] (0,0) -- (unitn);
  \end{scope}
  \node[note,anchor=north west,align=left] at (118,54)
    {共同下降区域\\$d_1>0$\\$d_2>d_1$};
  \node[label,anchor=south] at (151,44) {$\bm d_{\mathrm N}$};
  \draw[leader] (unitn) -- (151,44);
  \node[label,anchor=west] at (159,38) {$\bm d_{\mathrm M}$};
  \draw[leader] (unitm) -- (159,38);
  \node[label,anchor=east] at (133,12) {单位圆};
  \draw[leader] (134,13) -- (138.1,15.8);
  \node[note] at (142.5,4) {实际位移：$-\eta\bm d$};
\end{tikzpicture}

  \caption{固定原梯度的二维教学对照，各面板采用相同坐标尺度，$\bm g_1=(1,0)^{\top}$、$\bm g_2=(-2,2)^{\top}$。(a)虚线从原梯度端点连到PCGrad投影端点，两个投影向量求和为$\bm d_{\mathrm{PC}}$。(b)原梯度端点之间的线段是凸包，距原点最近的是$\bm d^\star=(4/13,6/13)^{\top}$。(c)点线圆限制单位方向，浅色扇区示意$d_1>0,d_2>d_1$的共同一阶下降组合，径向虚线边界不属于严格下降区域。实线和虚线分别是单位MGDA方向$\bm d_{\mathrm M}$与Nash方向$\bm d_{\mathrm N}$。各箭头均为梯度组合，参数位移需反向取$-\eta\bm d$；图不表示最终任务性能次序。}
  \label{fig:ket-sharing-gradient-geometry}
\end{figure*}
\FloatBarrier

零参考效用和同一可行集下，对数和的最优性还给出比例公平：对其他可行效用$u_k$，以最优效用$u_k^\star$归一化后的变化和$\sum_k(u_k-u_k^\star)/u_k^\star$不为正。把某项损失乘正常数只给对数目标加常量，故不改变这里的最优单位方向。这是局部效用的尺度性质，不替代任务优先级，也不意味着人群服务公平。群体成本须另设评价单位或约束，不能以总损失下降代替各任务、各群体收益。

专用协调器还需接受简单联合训练的实证对照。Xin等以英语到法语为一个任务，分别搭配英语到汉语、德语或罗马尼亚语，比较PCGrad、MGDA、GradNorm，以及其研究中的IMTL-G和随机损失加权对照。该研究以BLEU衡量译文与参考译文的片段匹配，分数越高越好；在其Transformer、Adam及所用调参范围内，专用方法的最终分数取舍没有超出简单采样比例扫描所得取舍，评价偏好改变还会改变方法排名\scite{ket-xin2022-mto}。该负面证据限于这些设置，没有测试Nash-MTL，不能据此否定所有协调方法。

上述更新处理的是参与训练的任务。检验从未共同训练的新任务是否受益，还须像第~\ref{sec:ket-transfer-task-splits}节那样留出任务并规定适应证据，不能只留出已见任务的图像。若冲突长期持续，也应检查分组、共享深度和输出规则，而非只调更新方向：当前结构没有保留的信息，优化器无法在这个结构内恢复。

\section{联合训练资源的分配}
\label{sec:ket-sharing-resources}

即使写下同一个联合目标，训练取得的证据和执行费用也可能不同。补标病害增加监督，提高病害抽取概率增加使用次数，扩大专家增加计算与表示自由度，三者不是同一种投入。下面分别核算样本标签、更新机会和计算存储，使实际训练与已定要求对应；总步数本身不足以完成这项比较。

\subsection{样本与标签资源的分配}
\label{sec:ket-sharing-labels}

有限标注可以用来覆盖更多植株，也可以给已有图像补齐其他任务。前者增加独立对象覆盖，后者让类别、病害与位置在同一对象上对应，便于学习输出关系。各任务均分、优先补薄弱任务、按对象成套标注，分别改变标签分布和可用对应；相同标签总数不能保证提供相同信息。

标签不全时，先只计算实际观察到的监督。以$m_{i,k}\in\{0,1\}$指示样本$i$有无任务$k$的可用标签，$n_{B,k}=\sum_{i\in B}m_{i,k}$为批次$B$中的可见数。只在$n_{B,k}>0$时对这些标签取平均：
\begin{equation}
  L_k(B)=
  \frac{1}{n_{B,k}}
  \sum_{i\in B}m_{i,k}
  \ell_k(f_k(\bm x_i),y_{i,k}).
  \label{eq:ket-sharing-label-mask}
\end{equation}
源码中的乘零是掩码记法，程序仍只能读取$m_{i,k}=1$的标签；缺标签时$y_{i,k}$不能被当作有效输入。若$n_{B,k}=0$，这个平均没有定义，也没有该任务的本批监督。未知病害状态不是健康，未知类别也不是负类。

例如，教学批次中图像1只有类别标签，图像2三种标签齐全，图像3只有病害与定位。各任务的两项可见损失分别为类别$(0.2,0.6)$、病害$(0.8,0.4)$、定位$(0.3,0.5)$，平均为$(0.4,0.6,0.4)$。若都除以完整批次大小3，就将各项乘以可见率$2/3$。可见数不同时，这种乘数还不同，因而会改变任务相对作用；两种平均均可作为目标约定，但不能无声互换。

删除缺失项并不消除选择偏差。专家若优先标明显病害，可见标签来自条件总体$m_k=1$，可见平均反映$\mathbb E[\ell_k\mid m_k=1]$，而非一般意义的总体风险。要恢复式~\eqref{eq:ket-sharing-task-risk}的对象，暂固定任务与预测器，假设$B$从该任务目标输入总体独立同分布抽取，再由标注机制揭示部分答案。

记固定预测器的样本损失为$\ell_{i,k}$。若$m_{i,k}\perp y_{i,k}\mid\bm x_i$，给定输入后可见性不再依赖未知答案；又若可见概率$e_k(\bm x)=\mathbb P(m_{i,k}=1\mid\bm x_i=\bm x)$已知且在目标支撑上处处为正，就可用逆可见概率补偿少被观察的输入。损失可积且批次非空时，采用完整批次作为分母：
\begin{equation}
  \widehat R_k(B)=\frac{1}{|B|}
    \sum_{i\in B}\frac{m_{i,k}}{e_k(\bm x_i)}\ell_{i,k}.
  \label{eq:ket-sharing-label-ipw}
\end{equation}
这一估计为何对应总体，可以先固定输入和答案，仅对标注是否出现取条件期望：
\begin{align}
  \mathbb E\!\left[
    \frac{m_{i,k}\ell_{i,k}}{e_k(\bm x_i)}
    \,\middle|\,\bm x_i,y_{i,k}\right]
  &=\frac{\ell_{i,k}}{e_k(\bm x_i)}
    \mathbb P(m_{i,k}=1\mid\bm x_i,y_{i,k})
  \nonumber\\
  &=\ell_{i,k}.
  \label{eq:ket-sharing-label-ipw-expectation}
\end{align}
第二个等号使用给定输入后的缺失独立性，随后再对输入答案取期望，就得到$\mathbb E[\widehat R_k(B)]=R_k(f_k)$。未见损失在这里只是理论随机变量，程序仍只累加可见项。若代入估计概率、可见性依赖未观察答案，或某区域概率为零，严格无偏所需依据便缺失，不能照用这个结论。

校正以后分母也必须对应原总体。假设教学批次有四项损失，均为1，可见概率均为$1/2$，恰观察到两个标签，则加权和4除以完整大小4得到1，除以可见数2却得到2。另可选择权重和作分母，称为\term{自归一化估计}（self-normalized estimator），即按本批实际权重总量重定平均尺度：
\begin{equation}
  \widehat R_{k,\mathrm{SN}}(B)
  =\frac{\sum_{i\in B}m_{i,k}\ell_{i,k}/e_k(\bm x_i)}
    {\sum_{i\in B}m_{i,k}/e_k(\bm x_i)}.
  \label{eq:ket-sharing-label-self-normalized}
\end{equation}
分母为正时才能计算这个比值。随机分母使其有限样本下一般有偏；常损失构造中同样得到1，只是该例的特殊结果，不说明它与完整批次估计恒等。

没有可见标签的批次能检验实现是否混淆了这几种对象。完整批次估计的值为0，表示本次随机贡献为空，不表示真实损失为0；可见平均和自归一化比值则未定义。监督反传可以省略，但风险平均仍保留完整批次计数，预定任务权重也不因可见任务少而重新归一化。固定模型时，零贡献批次被删去后，条件平均变为$R_k(f_k)/\mathbb P(n_{B,k}>0)$，已经改变了风险口径。

用于调权的学习进度不是这个随机贡献。GradNorm缺标时应保留最近有效进度；一种保守的教学处理是本步暂缓整组调权、沿用原权重，监督梯度仍按选定估计式计算。异步更新或只对可见任务调权须另定规则，不能把随机零值记成任务已学好。风险估计无偏还要在可交换期望与求导时才给出无偏梯度，更不能直接推到依赖随机状态的动态调权更新。相关校正思路见第~\ref{sec:ket-update-weighting}节。

样本费用应按实际证据计数。三种标签对应三项监督，却不把同一图像变成三张独立照片；原图增强和同株连续拍摄也应与独立对象区分。按植株划分测试时，同株记录须共同隔离。标签查询与核验属于知识获取，本节安排已有标签的使用，重复抽取不计为新增证据。

将任务改成自然语言接口时，也要保留这些单位。T0用多个模板描述数据集任务，并以留出任务类别检查未参与共同训练的要求\scite{ket-sanh2022-t0}。未见问法、未见数据来源和未见预测类别分别改变措辞、数据集和任务范围；一种留出不能代替另外两种。报告任务泛化前，应明确留出了哪一层。

FLAN的扩展研究混合多种指令数据，并讨论任务、模板与来源重复的核算\scite{ket-chung2024-flan}。同样，把“是否患病”改写多次不应计成多项独立能力。合并时分别检查对象、改写输入和答案，防止重复内容跨入评价集，并说明增加的是任务要求、独立观测还是表达形式。

\subsection{训练更新机会的分配}
\label{sec:ket-sharing-sampling}

即使标签清单不变，不同抽取次数也会改变任务作用。先考察采样前权重固定的线性加权更新，便于算清概率与权重如何相乘。PCGrad、MGDA和Nash-MTL需要同一状态下的多项梯度；分步收集或复用旧梯度属于额外近似，不能直接套用下面线性期望的结论。

每步先以$p_k>0$抽任务$k$，再从该任务的已标注训练集均匀抽样，用$w_k$加权所取得的梯度。固定当前参数、概率和权重，对任务选择及任务内样本取期望，得到
\begin{equation}
  \mathbb E[\widehat{\bm g}]
  =\sum_{k=1}^{K}p_kw_k
    \nabla\hat R_k.
  \label{eq:ket-sharing-effective-objective}
\end{equation}
所以当前期望更新对应$\sum_kp_kw_k\hat R_k$，不只是写出的损失权重。若任务内部又按难度偏向抽样，条件样本分布也会改变，须继续校正或改写目标。

令$n_k$为任务$k$的可抽标签记录数。取$p_k=n_k/\sum_jn_j$时，等权损失对应全部标签记录的样本平均；取$p_k=1/K$时，对应各任务平均。两者是不同要求。总步数$T$下抽取次数的期望为$Tp_k$，实际运行会波动。增频让任务更多次产生梯度，增权只放大已抽到的梯度；两者都不会自行增加独立标签。

若要固定目标而另选采样频率，先以$q_k\geq0,\sum_kq_k=1$指定目标系数；这里$q_k$不再是GradNorm中的损失进度比。抽到任务$k$时用$q_k/p_k$补偿其频率，就有
\begin{equation}
  \widehat{\bm g}
  =\frac{q_k}{p_k}\nabla\ell_k,
  \qquad
  \mathbb E[\widehat{\bm g}]
  =\sum_kq_k\nabla\hat R_k.
  \label{eq:ket-sharing-importance-sampling}
\end{equation}
$q_k>0$的任务必须有正抽取概率，且本次估计把$q_k/p_k$视为常量，不对采样决策求导。概率可根据训练历史调整，只要给定历史后仍满足条件，便保留对应条件期望。另有损失权重时，应核对最终系数乘积，避免重复补偿。

相同期望目标仍可能有不同波动。设$\bm g_k$为任务内随机样本梯度，$\bm\mu_k=\mathbb E[\bm g_k]$，各项二阶矩有限。展开校正梯度相对均值的平方范数，有
\begin{align}
  &\mathbb E\left\|
    \widehat{\bm g}-\sum_kq_k\bm\mu_k
  \right\|_2^2
  \nonumber\\
  &\quad=
  \sum_k\frac{q_k^2}{p_k}\mathbb E\|\bm g_k\|_2^2
  -\left\|\sum_kq_k\bm\mu_k\right\|_2^2.
  \label{eq:ket-sharing-sampling-variance}
\end{align}
二阶矩项含$1/p_k$，所以概率极小的任务偶尔被抽到时会产生很大校正作用。调整概率需要看所有项的二阶矩及单步费用，并不只是盲目增频；小样本池增频仍是在重复证据。概率下界可限制校正幅度，截断权重也可减少极端更新，但后者一般引入偏差，应记录改变。

静态分配以外，可以预定早期偏类别、后期增加病害和定位的课程，也可以按验证集未达标准的任务反馈调整。比较要固定总更新或总计算预算，记录各阶段实际次数与校正方式，最终测试不参与反馈选择。未校正的动态概率还会动态改变有效目标，不能把这种变化只解释成求解同一目标更高效。

\subsection{计算与存储预算的分配}
\label{sec:ket-sharing-budget}

同一图像同时请求类别、病害和位置时，共同骨干可以只算一次，但不同图像不能复用同一次激活。存储则另算：假设骨干参数量为$P_{\mathrm b}$，每个头均为$P_{\mathrm h}$，三个独立模型需$3(P_{\mathrm b}+P_{\mathrm h})$，共享骨干及三个头需$P_{\mathrm b}+3P_{\mathrm h}$。省下的是重复骨干参数，不代表所有调用都节省同样计算；只服务类别时还要解释额外任务训练费用的用途。

取教学参数量$P_{\mathrm b}=12$、$P_{\mathrm h}=1$，单位为百万，分别训练保存39，共享保存15，单任务共享路径涉及13。这个账目没有比较预测质量。若都允许总参数39，共享骨干可扩为36加三个头，或留出私有分支，但还需实测；相同参数数目也不保证浮点运算、时延或训练显存相同。

条件专家还有两种参数计数。\term{总容量}（total capacity）指全部保存参数，\term{激活容量}（active capacity）指一次请求实际参与计算的参数。设前端为$P_0$，每个专家为$P_{\mathrm e}$，任务门控和头分别为$P_{\mathrm g,k},P_{\mathrm h,k}$。以$A_k(\bm x)$表示输入上实际执行的专家集合，两个计数为
\begin{align}
  P_{\mathrm{total}}
    &=P_0+EP_{\mathrm e}
      +\sum_k(P_{\mathrm g,k}+P_{\mathrm h,k}),
  \nonumber\\
  P_{\mathrm{active},k}
    &=P_0+|A_k(\bm x)|P_{\mathrm e}
      +P_{\mathrm g,k}+P_{\mathrm h,k},
  \label{eq:ket-sharing-expert-budget}
\end{align}
稠密MMoE通常执行全部$E$个专家，只有增加稀疏机制才可能减少$|A_k|$。多个任务读同一图像时，专家若可缓存，其执行范围是$\bigcup_kA_k(\bm x)$，不能直接按任务相加。未执行参数仍可能占设备存储，所以激活少不等于驻留少。

以百万参数作教学计数，取$P_0=4,E=4,P_{\mathrm e}=6$，三个门控各$0.1$、头各1，总量为31.3。单任务稠密路径为29.1，另加只执行一个专家的稀疏机制后路径为11.1，但仍需保存31.3。实际费用还包含门控、输入聚合、批次形状及跨设备传输，参数计数不能直接换成时延。

训练显存还要加上梯度与优化器状态。若$P$个可训练参数均用32位浮点格式FP32，每个数占4字节，Adam分别保存参数、梯度及两个动量，四份数组共约$16P$字节，不含激活和缓冲。15百万参数为240\,MB，单份参数为60\,MB，MB按$10^6$字节计。部分专家本步不活跃，也通常仍保留已训练参数的优化器状态；混合精度副本、重计算、分片或卸载会改变账目。批次大小、分辨率和中间激活影响峰值，执行与通信细节见\BookPartRef{part:systems}{卷12《系统》第1部分}。

表~\ref{tab:ket-sharing-resource-controls}固定类别、病害和定位三项教学数据，分别有$600,300,100$条标签记录；每步抽一任务一记录，共1000步，数据、优化器与损失单位不变。$C_0$表示基准共享容量，容量变化不预设计算线性变化。读表可见，改权重与改采样能得到相同的任务平均目标，却保留不同机会和方差；改容量则改变候选模型及单步费用。

\begin{table*}[htbp]
  \centering
  \small
  \begin{tabularx}{\linewidth}{p{18mm}p{36mm}p{25mm}p{28mm}X}
    \toprule
    教学方案 & 抽取概率$\bm p$及权重$\bm w$ & 期望目标系数$p_kw_k$ & 抽取次数的期望 & 模型容量与费用核算 \\
    \midrule
    基准 &
    $\bm p=(0.6,0.3,0.1)$\newline
    $\bm w=(1,1,1)$ &
    $(0.6,0.3,0.1)$ &
    $(600,300,100)$ &
    容量为$C_0$；优化标签记录平均，单步费用作基准。 \\
    只改权重 &
    $\bm p$同基准\newline
    $\bm w=(5/9,10/9,10/3)$ &
    $(1/3,1/3,1/3)$ &
    $(600,300,100)$ &
    容量仍为$C_0$；抽取次数不变，低频任务的梯度放大。 \\
    只改采样 &
    $\bm p=(1/3,1/3,1/3)$\newline
    $\bm w=(1,1,1)$ &
    $(1/3,1/3,1/3)$ &
    各任务$1000/3$ &
    容量仍为$C_0$；小样本任务重复增多，按各任务单步费用计总计算。 \\
    只改容量 &
    $\bm p,\bm w$同基准 &
    $(0.6,0.3,0.1)$ &
    $(600,300,100)$ &
    容量增为$2C_0$；抽取与目标不变，重新测量单步成本。 \\
    联合改变 &
    $\bm p=(0.4,0.3,0.3)$\newline
    $\bm w=(5/6,10/9,10/9)$ &
    $(1/3,1/3,1/3)$ &
    $(400,300,300)$ &
    容量为$2C_0$；权重补偿抽取概率，分别核算机会与扩容。 \\
    \bottomrule
  \end{tabularx}
  \caption{相同任务数据下改变权重、抽取概率与容量的教学对照，数字不含实测性能，向量顺序为类别、病害、定位。抽取次数表示随机训练期望，不保证单次运行恰好达到。每步更新共享块及被选任务的监督路径，其他私有头不因此获得直接监督。目标系数相同仍可有不同抽取频率；1000步也不表示任务单步费用或容量改变后总计算相同。}
  \label{tab:ket-sharing-resource-controls}
\end{table*}

资源账目应与逐任务收益并列。类别准确率和病害召回率越高越好，定位误差越低越好，不能把三者原值直接平均。先报告各项原始指标和最低标准，再在已定义效用或基线归一化下汇总任务平均、兼容指标的样本平均与最弱任务结果，避免均值掩盖退化。标签、独立对象、重复次数、训练计算、参数与训练状态存储，以及部署调用条件均属于投入；第~\ref{chap:ket-evaluation}章给出统一评价要求。

\begin{samepage}
\section{本章小结}
\label{sec:ket-sharing-summary}

花卉类别、病害与定位有共同线索，也有不能丢失的差异。能否共享，需要把交付责任、共享位置、更新准则和资源安排放在同一方案中检验：选哪些任务成组，决定它们在哪里交换信息，再确认损失与采样产生了预定更新，并为各项能力保留足够数据和容量。共同参数或统一接口只让交换成为可能，收益仍由同评价条件、约定预算下的单任务与分组对照确定。

\end{samepage}

只交付类别时，属性目标以类别收益和额外成本决定去留；三项都交付时，各项退化分别承担责任；地区副本需要个性化时，还须允许共同中心或专家之外的局部变化。训练负内积不是长期负迁移，参数负相关也不是受损指标。失败时可据此分别核查关系测量、信息保留、尺度与优先级、标签可见性和采样预算，不把全部原因归入一个权重。

一项设计改变会影响其他项的证据。重新分组需重测亲和性，扩展专家需重查使用分布与容量，更改抽样需重算有效目标；收益比较仍保持任务单位、数据边界和成本口径一致。本章检查的是已见任务共同形成经验的效果，新任务迁移要留出任务并规定适应权限，长期保留与总体证据强度则分别回到知识维护和统一评价。
